\section{Finite characteristic points and propagation across components}
\label{sec:arithmetic-component-propagation}

\paragraph{What this chapter proves.}
The arithmetic construction uses two modularity realizations on one
frame, one in rank two and one in rank three.  Their modularity ideals
belong to different rings.  Component propagation makes the rank-two
ideal vanish on every relevant component, uses that information in
rank three, and produces an ideal controlling every nontriangular
rank-three point.  The last step is the original theorem
\(\texttt{FinPtOf}\): that ideal forces the characteristic polynomial
of every permitted characteristic-\(\ell\) point to be the prescribed
residual polynomial.

We prove this last theorem in full, including its original
\(\texttt{IsCD}\) test rings.  We also prove the original local and
global component-propagation principles, their transfer through flat
local charts, and the one-dimensional prime-avoidance theorem used
to find the local propagation points.  These are substantive inputs
to the two rank-specific propagation theorems.  The separate
deformation-space dimension, connectedness and rank-two-to-rank-three
corner constructions are stated explicitly at the end; their
conclusions are not assumed to be already reconstructed here.

There is one useful simplification in the finite-point argument.
The source realizes a point whose reduction may be triangular.
The actual \(\texttt{FinPt}\) consumer assumes the stronger,
literal condition \(\texttt{IsLiftA}\), which includes diagonal
reduction.  For these points the elementary subring
\(k+\mathfrak m_A\) gives the required realization directly.
No arbitrary coefficient extension or new frame is chosen.

\paragraph{Notation and the supplied universal space.}
Let \(O\) be the original complete discrete valuation ring with
finite residue field \(k\) of characteristic \(\ell\), and let
\(F\) be an original frame satisfying \(F.\mathrm{Ok}(C_0)\).
Write \(G_0=F.G_0\), \(G_v=F.G_v\), \(I_v=F.I_v\), and
\(\Psi=F.\Psi\).  The local maps are
\(\operatorname{loc}_{0,v}:G_v\to G_0\); their two branches are
\(\operatorname{loc}_{2,(v,0)}=\operatorname{loc}_{0,v}\) and
\(\operatorname{loc}_{2,(v,1)}=\operatorname{conj}\circ
\operatorname{loc}_{0,v}\).  Put
\[
 w(2)=(1,-1),\qquad w(3)=(1,0,-1).
                                                        \tag{ACP.1}
\]
The frame involution has
\(g^*=\operatorname{conj}(g^{-1})\) and
\(\Psi(g^*)=\Psi(g)^{-1}\).  Let \(J_n\) be the matrix with
ones on the antidiagonal and zeros elsewhere.

An \(\mathrm{NRep}(F,n,C)\) consists of a representation
\(r:G_0\to\mathrm{GL}_n(C)\) and characters
\(\psi_{v,i}:G_v\to C^\times\), with
\[
 r(g^*)=J_nr(g)^tJ_n,
 \qquad r(F.g_0)-r((F.g_0)^*)\ \text{diagonal},
                                                        \tag{ACP.2}
\]
and the ordered identity
\[
 \prod_{i=0}^{n-1}
 \bigl(r(\operatorname{loc}_{0,v}s_i)
                         -\psi_{v,i}(s_i)I_n\bigr)=0
 \quad\text{for every }(s_0,\ldots,s_{n-1})\in G_v^n.
                                                        \tag{ACP.3}
\]
Products in (ACP.3) have the displayed order.
Adic continuity means that every matrix entry and every character
value is locally constant modulo each power of the maximal ideal.
The condition \(\mathrm{IsLiftA}(O,r,\psi)\) means this continuity,
the specified residual characters \(\Psi^{w(n)_i}\), and zero
residual off-diagonal entries.  In particular it says more than
equality of residual semisimplifications.

A supplied \(S:\mathrm{Space}(O,F,n)\) has a complete Noetherian
local ring \(R\) with residue field \(k\), and a universal
\((r^{\mathrm{univ}},\psi^{\mathrm{univ}})\) for these continuous
diagonal-residual lifts over complete Noetherian local
\(O\)-algebras with residue field \(k\).  Write
\[
 \begin{split}
 \mathcal B&=\overline{O[\text{weight-zero entry monomials},
                         \ \psi^{\mathrm{univ}}_{v,i}(s)]},\\
 \mathcal A&=\mathcal B[\,r^{\mathrm{univ}}(g)_{ij}\,]
                                          \ \subset R.
 \end{split}                                             \tag{ACP.4}
\]
The closure is the \(\mathfrak m_R\)-adic closure.  The weight of
an entry \((i,j)\) is \(w(n)_i-w(n)_j\).  Thus a monomial belongs
to the displayed weight-zero set precisely when its entry weights
sum to zero.  Let \(u\) denote the universal representation over
\(\mathcal A\), let \(\iota:\mathcal B\hookrightarrow\mathcal A\)
be the inclusion, and set
\[
 \begin{split}
 \mathfrak m_{\mathcal B}&=\mathcal B\cap\mathfrak m_R,
 \quad\pi(\mathfrak p)=\iota^{-1}(\mathfrak p),\\
 \lambda&=\mathfrak m_O\mathcal B+
  \bigl(\psi_{v,i}^{\mathrm{univ}}(s)-1:
       s\in I_v,\ \Psi(\operatorname{loc}_v s)=1\bigr).
 \end{split}                                             \tag{ACP.5}
\]
The exact predicate \(S.\mathrm{FinLam}(I)\) is
\(\mathfrak m_{\mathcal B}\subset\sqrt{I+\lambda}\).

Let \(\mathfrak t_+\), respectively \(\mathfrak t_-\), be the
ideals of \(\mathcal A\) generated by all entries below,
respectively above, the diagonal.  A prime \(\mathfrak p\) is
\emph{off triangular} if neither ideal is contained in
\(\mathfrak p\).  The original condition \(\mathrm{OffFin}(S)\)
asserts the existence of one ideal \(I\subset\mathcal B\) with
\[
 \mathfrak m_{\mathcal B}\subset\sqrt{I+\lambda},
 \qquad I\subset\pi(\mathfrak p)
 \quad\text{for every off-triangular prime }\mathfrak p.
                                                        \tag{ACP.6}
\]
This same ideal must work before the later test ring and point are
chosen.

We use elementary local algebra, the principal ideal theorem,
finite prime avoidance, the Krull intersection theorem for
Noetherian local rings, and flat going down.  The arguments that
apply these facts, and the relevant constructions, are given below.
The frame's local generator and uniform class-exponent fields are
the literal fields reconstructed in LFC3 and LFC6; they are not
replaced by new finiteness assumptions depending on a test point.

\subsection{Realizing the exact test points without a larger construction}

% BEGIN ACP1
\paragraph{ACP1. The complete ring with the original residue constants.}
\label{claim:ACP1}
Let \(A\) satisfy the original \(\mathrm{IsCD}(O,A)\), and
suppose the specified maps form a scalar tower
\(O\to k\overset{j}{\longrightarrow}A\).
Thus \(A\) is a complete discrete valuation domain, its residue
field \(k_A\) is finite, and the given \(O\)-map is local.
Then
\[
 C=j(k)+\mathfrak m_A
   =\{a\in A:\overline a\in\overline j(k)\subset k_A\}
                                                        \tag{ACP.7}
\]
is a complete Noetherian local \(O\)-algebra with residue field
\(k\).  Its maximal ideal is \(\mathfrak m_A\), viewed as a
subset of \(C\), and
\[
 \mathfrak m_C^N=C\cap\mathfrak m_A^N\quad(N\ge0),
 \qquad \mathfrak m_C^N=\mathfrak m_A^N\quad(N\ge1)
                                                        \tag{ACP.8}
\]
under the inclusion in \(A\).  The \(O\)-structure is the
specified one.  In particular \(C\) satisfies the original
\(\mathrm{IsCO}(O,C)\) contract.

\begin{proof}
A unital map from the field \(k\) into the domain \(A\) is
injective.  Its composite with \(A\to k_A\) is also injective:
a nonzero element of \(k\) has its inverse already in \(A\),
so cannot lie in \(\mathfrak m_A\).  We may identify \(k\)
with these two images.  Since \(k_A\) is a finite set,
\([k_A:k]\) is finite.  This is supplied by \(\mathrm{IsCD}\);
no finite-degree hypothesis has been added.

The inverse image description in (ACP.7) makes \(C\) a subring.
If \(c\in C\setminus\mathfrak m_A\), its inverse in \(A\)
has residue \(\overline c^{-1}\in k\), hence also lies in \(C\).
Therefore \(C\) is local and \(\mathfrak m_C=\mathfrak m_A\).
Its residue map has image exactly \(k\), and the map from \(O\)
induces its prescribed residue isomorphism.

Choose a uniformizer \(t\) of \(A\).  For \(N\ge1\), every
element \(t^Na\in\mathfrak m_A^N\) is the product of \(N-1\)
copies of \(t\in\mathfrak m_C\) and the element
\(ta\in\mathfrak m_C\).  This proves the nontrivial inclusion
in (ACP.8); the reverse inclusion is immediate.  The assertion
for \(N=0\) is its intersection version.  The subring \(C\) is
closed in \(A\), since it is the inverse image of a subset of the
discrete residue field.  Equation (ACP.8) identifies its adic
Cauchy sequences and their limits with those in this closed
subring.  Hence \(C\) is complete and separated.

Here is a direct Noetherian argument.  Evaluation at \(t\) embeds
\(k[[T]]\) in \(A\).  Indeed, the first nonzero coefficient of
a series gives its exact \(t\)-adic order, so evaluation is
injective.  Let \(b_1,\ldots,b_d\in A\) lift a \(k\)-basis of
\(k_A\).  Every \(a\in A\) has a convergent expansion
\[
 a=\sum_{n\ge0}t^n\sum_{i=1}^d c_{n,i}b_i,
 \qquad c_{n,i}\in k.
                                                        \tag{ACP.9}
\]
To construct it, match the residue of \(a\) by a linear
combination of the \(b_i\), subtract, divide by \(t\), and
repeat.  Completeness makes the remainders tend to zero.
Regrouping the finite number of basis vectors shows that \(A\)
is a finite module over \(k[[t]]\).  The ring \(C\) contains
\(k[[t]]\) and is a submodule of that finite module.  Since
\(k[[T]]\) is Noetherian, \(C\) is a finite Noetherian module
over it, and therefore is a Noetherian ring.  If the zero-maximal-
ideal case is allowed in a DVR convention, then \(A=k_A\) and
\(C=k\), and all these assertions are immediate.

Finally \(O\to C\) factors as its original residue map followed
by the specified embedding of \(k\).  It is local, and every
element of \(C\) differs from an element of \(O\) by an element
of \(\mathfrak m_C\).  Together with completeness and
Noetherianity these are exactly all the \(\mathrm{IsCO}\) fields.
\end{proof}
% END ACP1

% BEGIN ACP2
\paragraph{ACP2. Every original finite-point test is realized on the supplied space.}
\label{claim:ACP2}
Let \(n=2\) or \(3\), let \(S:\mathrm{Space}(O,F,n)\) be
supplied, and let \(A\) and its coefficient structures be as in
ACP1.  For every \((r,\psi):\mathrm{NRep}(F,n,A)\) satisfying
\(\mathrm{IsLiftA}(O,r,\psi)\), there is an \(O\)-algebra map
\[
 \xi:\mathcal A\longrightarrow A
 \quad\text{with}\quad \xi_*u=(r,\psi).
                                                        \tag{ACP.10}
\]
Moreover \(\mathfrak m_{\mathcal B}\mathcal A\) maps into
\(\mathfrak m_A\).  This is the exact realization needed for
the original \(\mathrm{FinPt}\) consumer, on its original space.

\begin{proof}
Use the explicitly constructed \(C=k+\mathfrak m_A\).
Every off-diagonal matrix entry lies in \(\mathfrak m_A\),
by the residual-diagonal field of \(\mathrm{IsLiftA}\).
Each diagonal entry has residue
\(\Psi(g)^{w(n)_i}\in k\), so lies in \(C\).
Each full ordinary character value has residue
\(\Psi(\operatorname{loc}_v s)^{w(n)_i}\in k\), so also lies
in \(C\).  Its inverse is the character value at \(s^{-1}\),
and lies in \(C\) as well.  Thus these are characters to
\(C^\times\), not merely functions with values in the subring.

The representation and all its normalization, polarization and
ordered-product identities restrict to \(C\), because
\(C\hookrightarrow A\) is injective.  Equation (ACP.8) shows
that local constancy modulo powers of \(\mathfrak m_A\)
is exactly the required local constancy modulo powers of
\(\mathfrak m_C\) for these functions.  Reduction to the
residue field of \(C\) is the original diagonal representation
over \(k\), and the ordinary characters have their original
residual values.  This is an \(\mathrm{IsLift}(O,-)\) over
the \(\mathrm{IsCO}\)-algebra \(C\).

The universal property included in the supplied \(S\) therefore
gives \(f:R\to C\) carrying its universal representation to
this restricted one.  Restrict \(f\) to \(\mathcal A\) and
include \(C\) in \(A\); this is \(\xi\).
No assertion of existence of \(S\) has been used here: it was
part of the input before the test ring was chosen.

For the last assertion, an \(O\)-algebra map between two
\(\mathrm{IsCO}\)-algebras preserves their residue maps.
To see this directly, let \(x\in\mathfrak m_R\).  If \(f(x)\)
had nonzero residue, choose \(a\in O^\times\) with that residue.
Then \(x-a\) would be a unit in \(R\), while
\(f(x)-a\in\mathfrak m_C\) would be a nonunit, a contradiction.
It follows that \(f(\mathfrak m_R)\subset\mathfrak m_C\).
The stated image condition follows from
\(\mathfrak m_{\mathcal B}\subset\mathfrak m_R\).
\end{proof}
% END ACP2

\subsection{The local polynomial identity and rigidity of diagonal characters}

% BEGIN ACP3
\paragraph{ACP3. An ordered product with a separator gives the ordered flag.}
\label{claim:ACP3}
Let \(C\) be any commutative local ring, \(H\) any group,
\(r:H\to\mathrm{GL}_n(C)\), and
\(\psi_i:H\to C^\times\) for \(0\le i<n\).
Suppose an element \(s_0\in H\) satisfies
\[
 \psi_i(s_0)-\psi_j(s_0)\in C^\times\quad(i\ne j),
 \qquad
 \det(X-\overline{r(s_0)})
       =\prod_i(X-\overline{\psi_i(s_0)}),
                                                        \tag{ACP.11}
\]
and suppose (ACP.3), with \(H\) in place of \(G_v\), holds
for every tuple.  There is one basis of \(C^n\) in which
every \(r(s)\) is upper triangular with ordered diagonal
\((\psi_0(s),\ldots,\psi_{n-1}(s))\).
This proves the full original \(\texttt{flag\_loc}\) statement;
no domain, completeness or continuity hypothesis is necessary.

\begin{proof}
Put \(T=r(s_0)\) and \(a_i=\psi_i(s_0)\).
The constant tuple in the ordered-product identity gives
\(\prod_i(T-a_i)=0\).  Define
\[
 E_i=\prod_{j\ne i}(T-a_j)(a_i-a_j)^{-1}.
                                                        \tag{ACP.12}
\]
Lagrange interpolation modulo the monic polynomial
\(\prod_i(X-a_i)\) gives
\(\sum E_i=1\), \(E_iE_j=0\) for \(i\ne j\),
\(E_i^2=E_i\), and \(TE_i=a_iE_i\).
These polynomial identities hold over \(C\), since the displayed
differences are units.  Modulo the maximal ideal, (ACP.11) has
distinct roots and each \(\overline E_i\) has a one-dimensional
image.  Choose lifts \(v_i\) of a basis of those images.
The vectors \(E_iv_i\) reduce to a basis, so the matrix having
them as columns has unit determinant.  They form a basis of
\(C^n\), in which the \(E_i\)'s are the coordinate projectors.

Fix \(s\in H\) and indices \(b\le a\).  In (ACP.3) put
\(s\) in position \(a\) and \(s_0\) in every other position.
Multiply on the left by \(E_a\) and on the right by \(E_b\).
Each factor before position \(a\) becomes the unit scalar
\(a_a-a_j\), and each factor after it becomes the unit scalar
\(a_b-a_j\).  Cancelling their product yields
\[
          E_a\bigl(r(s)-\psi_a(s)\bigr)E_b=0.
                                                        \tag{ACP.13}
\]
For \(b<a\), orthogonality gives \(E_ar(s)E_b=0\);
for \(b=a\), it gives \(E_ar(s)E_a=\psi_a(s)E_a\).
These are exactly the lower-zero and ordered-diagonal assertions
in the chosen basis, simultaneously for every \(s\).
For \(n=0\), the unique basis of the zero module gives the same
conclusion directly.
\end{proof}
% END ACP3

% BEGIN ACP4
\paragraph{ACP4. The literal local characteristic polynomial on the universal space.}
\label{claim:ACP4}
For every supplied \(S:\mathrm{Space}(O,F,n)\), with
\(n=2\) or \(3\), one has
\[
 \det\bigl(X-u(\operatorname{loc}_{0,v}s)\bigr)
       =\prod_i\bigl(X-\psi^{\mathrm{univ}}_{v,i}(s)\bigr)
                  \quad\text{in }\mathcal A[X]
                                                        \tag{ACP.14}
\]
for every \(v,s\).  This is the original \(\texttt{o1}\)
assertion, and it remains true after every algebra map out of
\(\mathcal A\).

\begin{proof}
The \(f3\) field of \(F.\mathrm{Ok}(C_0)\) gives
\(s_0\in G_v\) with
\(\Psi(\operatorname{loc}_v s_0)^2\ne1\).
For the two lists of weights (ACP.1), the corresponding powers
of this value are pairwise distinct.  Thus the differences of
\(\psi^{\mathrm{univ}}_{v,i}(s_0)\) are units of the local
ring \(R\).  The residual-diagonal field of the universal lift
gives (ACP.11), with \(C=R\).  Apply ACP3 to the restricted
universal representation and its exact local characters.
The characteristic polynomial of an upper-triangular matrix is
the product of its diagonal factors, proving (ACP.14) over \(R\).
All its coefficients belong to \(\mathcal A\), and
\(\mathcal A\hookrightarrow R\) is injective, so the identity
holds there.  Polynomial evaluation under any later algebra map
preserves the identity; no flatness of that map is needed.
\end{proof}
% END ACP4

% BEGIN ACP5
\paragraph{ACP5. Principal units have no torsion in a characteristic-prime local domain.}
\label{claim:ACP5}
Let \(C\) be a local domain of characteristic \(p>0\), with
\(p\) prime.  If \(x-1\in\mathfrak m_C\) and \(x^h=1\)
for an integer \(h>0\), then \(x=1\).
This is the full original \(\texttt{oneUnitTF}\) assertion.

\begin{proof}
Write \(h=p^am\) with \(p\nmid m\), and put \(y=x^{p^a}\).
Then \(y^m=1\) and \(y-1=(x-1)^{p^a}\in\mathfrak m_C\).
The geometric sum \(1+y+\cdots+y^{m-1}\) has residue \(m\),
which is a nonzero element of the prime field and hence a unit.
The identity
\[
          (y-1)(1+y+\cdots+y^{m-1})=y^m-1=0
\]
gives \(y=1\).  Thus \((x-1)^{p^a}=0\); the domain
hypothesis gives \(x=1\).
\end{proof}
% END ACP5

% BEGIN ACP6
\paragraph{ACP6. A bounded finite-character quotient determines the whole character.}
\label{claim:ACP6}
Let \(G\) be a topological group, \(T\subset G\), and \(h>0\).
Assume every homomorphism from \(G\) to a finite abelian
\(p\)-group, with open kernel and trivial on \(T\), has exponent
dividing \(h\).  Let \(C\) be a Noetherian local domain of
characteristic \(p\) with finite residue field.  Suppose
\(a,b:G\to C^\times\) satisfy:
\(a\) is adically continuous, \(\ker b\) is open,
\(a\equiv b\pmod{\mathfrak m_C}\), and \(a=b\) on \(T\).
Then \(a=b\) on \(G\).  Completeness of \(C\) is not required.

\begin{proof}
Put \(\chi=ab^{-1}\).  Its values lie in
\(1+\mathfrak m_C\).  Each ring
\(C/\mathfrak m_C^N\) is finite: its successive ideal
quotients are finite-dimensional vector spaces over the finite
residue field.  The image \(B_N\) of \(\chi\) in its unit
group is an abelian \(p\)-group.  Indeed, for
\(z\in\mathfrak m_C\), the characteristic-\(p\) identity
\((1+z)^{p^N}=1+z^{p^N}\) and \(p^N\ge N\) show that its
reduction is annihilated by \(p^N\).

The character \(\chi_N:G\to B_N\) has open kernel.
To verify this without any unstated topology on unit groups,
intersect the open kernel of \(b\) with a neighborhood of
\(1\) on which \(a\) is \(1\) modulo \(\mathfrak m_C^N\).
This intersection lies in \(\ker\chi_N\); a subgroup containing
an open identity neighborhood is open.  Moreover \(\chi_N\)
is trivial on \(T\).  By the uniform assumption,
\[
                  \chi(g)^h-1\in\mathfrak m_C^N
                  \quad\text{for every }N,g.
\]
Krull intersection in the Noetherian local domain gives
\(\chi(g)^h=1\).  ACP5, applied to \(\chi(g)\), gives
\(\chi(g)=1\) for every \(g\), as required.
\end{proof}
% END ACP6

% BEGIN ACP7
\paragraph{ACP7. The frame determines every multiplicative diagonal entry.}
\label{claim:ACP7}
Let \(C\) be a Noetherian local domain with finite residue field,
let \(j:k\to C\) be a ring map, and let
\((r,\psi):\mathrm{NRep}(F,3,C)\) be adically continuous.
Suppose, for every \(g,h\in G_0\) and \(i\),
\[
 \begin{split}
 r(gh)_{ii}&=r(g)_{ii}r(h)_{ii},\\
 r(g)_{ii}&\equiv j(\Psi(g)^{w_i})\pmod{\mathfrak m_C},\\
 s\in I_v,\ \Psi(\operatorname{loc}_v s)=1
       &\Longrightarrow r(\operatorname{loc}_{0,v}s)_{ii}=1.
 \end{split}                                             \tag{ACP.15}
\]
Then \(r(g)_{ii}=j(\Psi(g)^{w_i})\) for every \(g,i\).
This is the full diagonal-rigidity step used in
\(\texttt{FinPtOf}\).

\begin{proof}
The embedding \(j\) makes \(C\) a ring of characteristic
\(\ell\).  Fix \(i\).  Multiplicativity and the equation at
\(g^{-1}\) make \(a_i(g)=r(g)_{ii}\) a character to
\(C^\times\), with inverse \(r(g^{-1})_{ii}\).
The reference character is \(b_i(g)=j(\Psi(g)^{w_i})\).
Its kernel is open, since it contains the open kernel of the
frame character.  The first two requirements of ACP6 are now
the given continuity and congruence.

Choose the local detectors \(u_{v,a}\) from the exact
\(\mathrm{gen}\) field, and the exponent \(h>0\) from
\(\mathrm{cl}\).  These are fixed by the frame before \(C\)
or \(a_i\) is considered.  Put
\[
 T=\{\operatorname{loc}_{2,(v,b)}(u_{v,a}):
                        v\in V,\ b\in\{0,1\},\ a<d_v\}.
                                                        \tag{ACP.16}
\]
A finite abelian \(\ell\)-group character with open kernel
which kills \(T\) kills each full local inertia image:
apply \(\mathrm{gen}\) to its pullback along each continuous
\(\operatorname{loc}_{2,(v,b)}\).  The \(\mathrm{cl}\) field
then gives the same exponent \(h\) for that character.  Thus
the finite-character hypothesis of ACP6 holds for (ACP.16).

For branch \(b=0\), the third equation of (ACP.15) applies
to every detector because \(u_{v,a}\in I_v\) and its residual
frame character is \(1\).  For branch \(b=1\), use
\(\operatorname{conj}(g)=(g^{-1})^*\) in (ACP.2):
\[
 r(\operatorname{conj}g)_{ii}=r(g^{-1})_{2-i,2-i}.
\]
When all diagonal values at \(g\) are one, their inverse
values on the right are also one.  The reference character is
one on both branches, by the corresponding frame identity
for \(\Psi\).  Hence \(a_i=b_i\) on \(T\).
ACP6 gives equality everywhere, separately for each \(i\),
with no change in the exponent or detector choices.
\end{proof}
% END ACP7

\subsection{The complete finite-point theorem}

% BEGIN ACP8
\paragraph{ACP8. The five cycle products determine a three-by-three characteristic polynomial.}
\label{claim:ACP8}
For a matrix \(M=(m_{ij})\) over any commutative ring, suppose
\[
 m_{01}m_{10}=m_{02}m_{20}=m_{12}m_{21}=0,
 \qquad m_{01}m_{12}m_{20}=m_{02}m_{21}m_{10}=0.
                                                        \tag{ACP.17}
\]
Then \(\det(X-M)=\prod_{i=0}^2(X-m_{ii})\).

\begin{proof}
Expand the determinant over its six permutations.  The identity
permutation contributes \(\prod_i(X-m_{ii})\).  The three
transpositions contribute respectively
\[
 -(X-m_{22})m_{01}m_{10},\quad
 -(X-m_{11})m_{02}m_{20},\quad
 -(X-m_{00})m_{12}m_{21},
\]
and the two three-cycles contribute
\(-m_{01}m_{12}m_{20}\) and \(-m_{02}m_{21}m_{10}\).
They all vanish by (ACP.17).
\end{proof}
% END ACP8

% BEGIN ACP9
\paragraph{ACP9. Off-triangular test points collapse the weight-zero algebra.}
\label{claim:ACP9}
Let \(S:\mathrm{Space}(O,F,3)\) satisfy
\(\mathrm{OffFin}(S)\).  Let \(A\) have the original
\(\mathrm{IsCD}\) and scalar-tower structures, and let
\((r,\psi)\) be an \(\mathrm{IsLiftA}\) with
\[
 s\in I_v,\ \Psi(\operatorname{loc}_v s)=1
                         \Longrightarrow\psi_{v,i}(s)=1.
                                                        \tag{ACP.18}
\]
Choose \(\xi\) by ACP2.  If
\(\ker\xi\subset\mathcal A\) is off triangular, then
\[
 \xi(b)=j(\operatorname{res}_R b)\quad(b\in\mathcal B).
                                                        \tag{ACP.19}
\]
Consequently every weight-zero product of entries of \(r\)
which contains an off-diagonal entry is zero, and
\(r(g)_{ii}=j(\Psi(g)^{w_i})\) for every \(g,i\).

\begin{proof}
Write \(\varphi=\xi\circ\iota:\mathcal B\to A\).
Both \(\ker\xi\) and \(\ker\varphi\) are prime, because
\(A\) is a domain.  The ideal \(I\) from (ACP.6) lies in
\(\ker\varphi\).  The ideal \(\lambda\) also lies there:
its coefficient generators vanish because \(O\to A\) factors
through \(k\), and its ordinary-character generators vanish
by (ACP.18).  Hence
\[
 \mathfrak m_{\mathcal B}\subset\sqrt{I+\lambda}
                                      \subset\ker\varphi.
\]
Given \(b\in\mathcal B\), choose \(a\in O\) with residue
\(\operatorname{res}_R b\).  The element \(b-a\) lies in
\(\mathfrak m_{\mathcal B}\), so
\(\varphi(b)=\varphi(a)=j(\operatorname{res}_R b)\).
This proves (ACP.19), including the original coefficient map.

A weight-zero entry monomial belongs to \(\mathcal B\) by
definition (ACP.4).  If it contains an off-diagonal entry,
its universal residue is zero, since the universal lift is
diagonal on reduction.  Equation (ACP.19) makes its image zero.
A single diagonal entry is itself a weight-zero monomial, and
its residue is \(\Psi(g)^{w_i}\), proving the final assertion.
In particular the five products in (ACP.17) vanish: the weights
around each indicated cycle telescope to zero.
\end{proof}
% END ACP9

% BEGIN ACP10
\paragraph{ACP10. Triangular test points have the same exact diagonal characters.}
\label{claim:ACP10}
Keep the data and condition (ACP.18) of ACP9.  If \(\ker\xi\)
is not off triangular, then all the matrices \(r(g)\) are
upper triangular in the given basis, or all are lower triangular.
In either case
\[
                 r(g)_{ii}=j(\Psi(g)^{w_i})
                 \quad\text{for every }g,i.
                                                        \tag{ACP.20}
\]

\begin{proof}
Failure of the off-triangular condition says that
\(\mathfrak t_+\subset\ker\xi\) or
\(\mathfrak t_-\subset\ker\xi\).  These are assertions
about all generators of one fixed ideal, so they give one
triangular direction simultaneously for every \(g\).
Diagonal entries of products of matrices with this direction
multiply.  Thus the first line of (ACP.15) holds.
The second is the residual field of \(\mathrm{IsLiftA}\).

Take \(s\in I_v\) with \(\Psi(\operatorname{loc}_v s)=1\).
ACP4 and (ACP.18) give
\[
 \prod_{i=0}^2\bigl(X-r(\operatorname{loc}_{0,v}s)_{ii}\bigr)
                                                    =(X-1)^3.
\]
Evaluate at the \(i\)-th diagonal entry.  Its difference from
one has cube zero, hence is zero in the domain \(A\).
This proves the third line of (ACP.15).  Since a DVR is
Noetherian and \(\mathrm{IsCD}\) supplies a finite residue
field, ACP7 applies and proves (ACP.20).

Every product in (ACP.17) also vanishes in either triangular
direction.  Thus this case has exactly the same polynomial
conclusion as the off-triangular case, without requiring
\(I\subset\ker\varphi\) here.
\end{proof}
% END ACP10

% BEGIN ACP11
\paragraph{ACP11. Full original finite-point control and its fixed-frame application.}
\label{claim:ACP11}
For the original coefficient DVR \(O\), every
\(F.\mathrm{Ok}(C_0)\), and every supplied
\(S:\mathrm{Space}(O,F,3)\),
\[
               \mathrm{OffFin}(S)\ \Longrightarrow
                         \mathrm{FinPt}(O,F,3).
                                                        \tag{ACP.21}
\]
Explicitly, for every original \(\mathrm{IsCD}(O,A)\), every
specified scalar tower \(O\to k\to A\), and every
\(\mathrm{IsLiftA}\) satisfying (ACP.18),
\[
 \det(X-r(g))=
   \prod_{i=0}^2\bigl(X-j(\Psi(g)^{w_i})\bigr)
                            \quad(g\in G_0).
                                                        \tag{ACP.22}
\]
There is no added characteristic-zero, coefficient-flatness,
finite-module, or finite-extension-of-fraction-fields assumption.
Neither \(\mathrm{Big}\) nor an automorphy premise is needed
for this implication.  This proves the full original
\(\texttt{FinPtOf}\) contract.

\begin{proof}
Fix an arbitrary test ring and point under the displayed
quantifiers.  ACP1--2 realize it on the already supplied space.
Apply ACP9 or ACP10 according as its prime kernel is off
triangular or not.  In both cases the diagonal entries are
the exact specified characters and the five products (ACP.17)
vanish.  ACP8 therefore gives (ACP.22) for every \(g\).
The choice of the test point was arbitrary, so the original
universal quantifiers have all been retained.

For the later common-frame construction, keep the frame,
link and marked \(g_1\) supplied by AFE15 or AFE16.
Once a space \(S\) on that frame has been constructed with
\(\mathrm{OffFin}(S)\), (ACP.21) supplies the precise premise
of AFE20.  That already proved transport chooses
\(U=G_L\) before every later ordinary point and returns
\(\mathrm{FinEnd}\) on the same \(g_1\).  This application
does not assert that \(S\) or \(\mathrm{OffFin}(S)\) has been
constructed merely by stating them.  It identifies the exact
consumer of the finite-point theorem just proved.
\end{proof}
% END ACP11

\subsection{One-dimensional primes that avoid every forbidden locus}

% BEGIN ACP12
\paragraph{ACP12. A one-dimensional specialization avoiding a prescribed element.}
\label{claim:ACP12}
Let \((R,\mathfrak m)\) be a Noetherian local ring,
\(Q\subset R\) a prime ideal, and \(f\in\mathfrak m\setminus Q\).
There is a prime \(P\supset Q\) such that
\[
                    f\notin P,\qquad \dim R/P=1.
                                                        \tag{ACP.23}
\]

\begin{proof}
By the ascending-chain condition choose a prime \(P\) maximal
among primes containing \(Q\) and avoiding \(f\).
The local domain \(B=R/P\) has a nonzero nonunit
\(\bar f\), and every nonzero prime of \(B\) contains
\(\bar f\): otherwise its inverse image would contradict
the maximality of \(P\).

We prove directly that this forces \(\dim B=1\).
Let \(P_1,\ldots,P_s\) be the finitely many minimal primes
over \((\bar f)\).  They have height at most one by the
principal ideal theorem and are nonzero because \(B\) is a
domain and \(\bar f\ne0\).
Given \(0\ne x\in\mathfrak m_B\), choose a prime \(Q_x\)
minimal over \((x)\) and contained in \(\mathfrak m_B\).
It is nonzero and has height at most one.  By the preceding
property it contains \(\bar f\), and hence contains some
\(P_i\).  Strict containment would give the chain
\((0)\subsetneq P_i\subsetneq Q_x\), contradicting that
height bound.  Therefore \(Q_x=P_i\), so \(x\in P_i\).
Zero belongs to the same finite union.  We have proved
\(\mathfrak m_B\subset\bigcup_iP_i\).
Finite prime avoidance gives \(\mathfrak m_B\subset P_i\)
for one index, and thus \(\mathfrak m_B=P_i\).
Its height is at most one and at least one, since it contains
the nonzero element \(\bar f\).  A local ring has dimension
equal to the height of its maximal ideal, so \(\dim B=1\).
\end{proof}
% END ACP12

% BEGIN ACP13
\paragraph{ACP13. The full finite-family nice-prime existence theorem.}
\label{claim:ACP13}
Let \((R,\mathfrak m)\) be any Noetherian local ring,
\(\varpi\in\mathfrak m\), and \(N\ge1\).
Let \((\mathfrak a_i)_{i\in E}\) be a finite family of ideals
with \(\dim R/\mathfrak a_i<N\), and let \(I\) be an ideal
with \(\dim R/I\ge N+1\).  Then there is a prime \(P\)
such that
\[
 I\subset P,\qquad \varpi\in P,\qquad \dim R/P=1,
 \qquad \mathfrak a_i\not\subset P\quad(i\in E).
                                                        \tag{ACP.24}
\]
This proves the original \(\mathrm{NicePrimeEx}\), including
arbitrary \(I\), arbitrary \(\varpi\in\mathfrak m\), and an
arbitrary finite indexing type.  The family may be empty.

\begin{proof}
First
\(\dim R/(I,\varpi)\ge N\).
Here the familiar one-equation dimension bound requires no
catenarity assumption.  In a Noetherian local ring \(T\),
lift a system of parameters of \(T/(x)\), of length
\(d=\dim T/(x)\).  Together with \(x\), these generate an
ideal with maximal-ideal radical.  The generalized principal
ideal theorem gives \(\dim T\le d+1\).
Apply this to \(T=R/I\); the quotients are nonzero since
\(I\) and \(\varpi\) lie in the maximal ideal.

Choose a minimal prime \(Q\) over \((I,\varpi)\) with
\(\dim R/Q\ge N\).  Such a prime exists because there are
finitely many minimal primes and the dimension of a Noetherian
quotient is the maximum of the dimensions of its irreducible
components.  No \(\mathfrak a_i\) is contained in \(Q\):
that containment would imply
\(\dim R/\mathfrak a_i\ge\dim R/Q\ge N\).
Choose \(a_i\in\mathfrak a_i\setminus Q\).
As \(N\ge1\), the prime \(Q\) is not maximal; choose
\(g\in\mathfrak m\setminus Q\), and put
\[
                           f=g\prod_{i\in E}a_i.
\]
It belongs to \(\mathfrak m\) and avoids the prime \(Q\).
ACP12 gives a prime \(P\supset Q\), of quotient dimension
one, avoiding \(f\).  It contains \(I\) and \(\varpi\).
If \(\mathfrak a_i\subset P\), then \(a_i\in P\) and
\(f\in P\), a contradiction.  For an empty family the
product is one and exactly the same argument applies.
\end{proof}
% END ACP13

\subsection{Propagation inside a local deformation chart}

\paragraph{The exact local propagation data.}
Let \(R\) be a Noetherian ring and use the order of ideal
inclusion on \(\operatorname{Spec}R\).  For an ideal \(I\),
write \(\operatorname{qdim}(I)=\dim(R/I)\); the zero quotient
has dimension \(-\infty\), as in the source.  Fix an integer
\(D\ge0\), subsets
\(\mathcal T,\mathcal E,\mathcal G,\mathcal N,\mathcal H\)
of \(\operatorname{Spec}R\), and an ideal \(\mathfrak s\).
They denote the type-I, exceptional-type, good, nice, and thick
primes.  A \emph{component prime} is a minimal prime over
\(\mathfrak s\).  Two such primes \(P,Q\) are adjacent if
\[
                         \dim R/(P+Q)\ge D.
                                                        \tag{ACP.25}
\]
Two nice exceptional primes \(p,q\in\mathcal N\cap\mathcal E\)
are \emph{linked} when there is a thick component prime
\(e\in\mathcal H\) contained in both.  A chain always means
a finite chain, and a chain of length zero is permitted.

The original \(\mathrm{LocData.Ok}\) consists of the following
ten requirements.  We spell them out to distinguish the
propagation theorem from its later geometric inputs.
\begin{enumerate}
\item Good primes are upward closed.
\item Type-I primes are downward closed.
\item Every nice prime is type I or exceptional.
\item Every good prime \(I\) with \(\operatorname{qdim}(I)\ge D\)
      is contained in a nice prime.
\item If \(e\) is a thick component prime and
      \(e\subset I\) with \(\operatorname{qdim}(I)\ge D\),
      then \(I\) is contained in a nice exceptional prime.
\item Every type-I minimal prime of \(R\) is a component prime.
\item Any two component primes are connected by adjacency steps
      (ACP.25).
\item Every component prime is type I or thick.
\item If a nice type-I prime is good, all primes below it are good.
\item If a nice exceptional prime \(p\) has a good type-I prime
      below it, then along every chain of links from \(p\) to
      \(q\), every type-I prime below \(q\) is good.
\end{enumerate}
The last two are the literal \(P1\) and
\(\mathrm{flipTwoThree}\) properties.  They are inputs to
this local theorem; the theorem proves how they propagate across
all component primes.

% BEGIN ACP14
\paragraph{ACP14. The full local component-propagation theorem.}
\label{claim:ACP14}
Under exactly the ten conditions above, if \(I\) is good,
type I, and \(\operatorname{qdim}(I)\ge D\), then every
type-I minimal prime of \(R\) is good.  This is the full
original \(\mathrm{LocProp}\), with no local-ring hypothesis
on \(R\) and no extra disjointness condition on the five sets.

\begin{proof}
We first record a consequence that applies at any nice prime
\(p\).  If there is a good type-I prime below \(p\), then
every type-I prime below \(p\) is good.  If \(p\) is type I,
upward closure first makes \(p\) good and condition 9 applies.
If it is exceptional, apply condition 10 to the chain of
length zero.  We will call this the propagation property at
\(p\).

We propagate the following precise information along component
primes.  A component prime \(Q\) is in one of two states:
\begin{itemize}
\item \(Q\) is type I and good;
\item \(Q\) is thick, and there is a nice exceptional prime
      \(p_0\) with a good type-I prime below it, together with
      a link chain from \(p_0\) to a nice exceptional prime
      \(p\) containing \(Q\).
\end{itemize}
In either state, if \(Q\) is also type I, it is good.
For the second state this follows from condition 10.

Here is a useful extension of either state.  Suppose
\(Q\subset M\) and \(\operatorname{qdim}(M)\ge D\).
Then there is a nice prime \(p'\supset M\) such that every
type-I prime below \(p'\) is good.  In the first state,
\(M\) is good by upward closure; condition 4 gives \(p'\),
and the propagation property at \(p'\) applies to \(Q\).
In the second state, condition 5 applied to the thick
component \(Q\) gives a nice exceptional \(p'\supset M\).
The primes \(p,p'\) are linked by \(Q\).  Append this link
to the existing chain and apply condition 10.

Now suppose \(Q,Q'\) are adjacent and the stated information
is available at \(Q\).  There is a prime \(M\supset Q+Q'\)
with \(\operatorname{qdim}(M)\ge D\): take the first prime
in a chain witnessing the dimension bound for \(R/(Q+Q')\).
The preceding extension gives a nice \(p'\supset M\) at
which every type-I prime below it is good.
If \(Q'\) is type I, it is therefore good, giving the first
state.  Otherwise condition 8 makes \(Q'\) thick.
Condition 5, applied to \(Q'\subset M\), gives a nice
exceptional prime \(q\supset M\).

If \(Q\) was in the first state, the good type-I prime
\(Q\subset q\) makes \(q\) itself a starting prime with
the required witness.  Use the length-zero chain at \(q\)
to put \(Q'\) in the second state.
If \(Q\) was in the second state, its previous ending prime
\(p\) and the new \(q\) are linked by the thick component
\(Q\), because both contain \(Q\).  Appending this link
puts \(Q'\) in the second state.  This proves the induction
step along every adjacency chain.  Notice that no assertion
that a thick prime is not type I was used.

To begin the induction, choose a minimal prime \(P_0\subset I\).
Condition 2 makes it type I, and condition 6 makes it a
component prime.  Condition 4 gives a nice \(p\supset I\).
The propagation property at \(p\), using the given good
type-I prime \(I\), makes \(P_0\) good.  Thus \(P_0\) is in
the first state.  Any type-I minimal prime \(P\) of \(R\)
is also a component prime, and condition 7 gives a finite
adjacency chain from \(P_0\) to \(P\).  The induction makes
\(P\) good, proving the assertion.
\end{proof}
% END ACP14

\subsection{From local component propagation to the whole spectrum}

\paragraph{The exact global propagation data.}
Let \(X\) be a Noetherian topological space, with subsets
\(T,G,C,Z,E\).  They are the type-I locus, good locus, core,
chosen closed-point locus, and exceptional locus, respectively.
Call an irreducible component \(Y\) relevant if \(Y\cap T\ne\varnothing\).
The original \(\mathrm{GlobData.Ok}\) asks that:
\begin{enumerate}
\item \(Z\subset C\);
\item \(C\) is preconnected;
\item every closed set meeting \(C\) meets \(Z\);
\item \(C\subset\overline{Z\setminus E}\);
\item every relevant component meets \(C\);
\item whenever two relevant components meet at a point of \(Z\),
      goodness of the whole first component implies goodness of
      the whole second;
\item every point of \(Z\setminus E\) belongs to a relevant component;
\item at least one relevant component is contained in \(G\).
\end{enumerate}
Neither openness of \(T\) nor closedness of \(G\) is assumed
in this abstract global theorem.  Those additional properties
will enter the chart-transfer theorem separately.

% BEGIN ACP15
\paragraph{ACP15. The full global component-propagation theorem.}
\label{claim:ACP15}
Under exactly the eight global conditions above, every relevant
component is contained in \(G\); in particular \(T\subset G\).
This is the full original \(\mathrm{GlobProp}\).

\begin{proof}
A Noetherian topological space has finitely many irreducible
components.  Let \(G_0\) be the union of its relevant
components that are contained in \(G\), and let \(B_0\) be
the union of the remaining relevant components.  Both are
closed, being finite unions of closed irreducible components.
Condition 7 places \(Z\setminus E\) in \(G_0\cup B_0\),
so condition 4 places the whole core \(C\) in this closed
union.  Condition 8, together with condition 5, shows that
\(C\cap G_0\ne\varnothing\).

Suppose a bad relevant component existed.  By condition 5
it would also give \(C\cap B_0\ne\varnothing\).
Preconnectedness of \(C\), applied to the two closed sets
covering it, gives a point of \(C\cap G_0\cap B_0\).
Condition 3, applied to the closed set \(G_0\cap B_0\),
therefore gives a point \(z\in Z\cap G_0\cap B_0\).
Choose a good relevant component \(Y\) and a bad relevant
component \(Y'\) containing \(z\).  Condition 6 makes
\(Y'\subset G\), contradicting its definition.
Thus no bad relevant component exists.

Every point of \(T\) belongs to an irreducible component,
which is automatically relevant.  The proved component
assertion therefore gives \(T\subset G\).
\end{proof}
% END ACP15

% BEGIN ACP16
\paragraph{ACP16. Flat local maps give the full required component chart.}
\label{claim:ACP16}
Let \(j:\operatorname{Spec}C\to X\) be an open embedding
into an arbitrary topological space.  Let \(R\) be a
commutative local \(C\)-algebra, flat over \(C\), and put
\[
 f=j\circ\operatorname{Spec}(C\to R),
 \qquad x=f(\mathfrak m_R).
\]
Then \(f\) is continuous, sends the closed point to \(x\),
and has the following two properties:
\begin{enumerate}
\item for every minimal prime \(P\) of \(R\),
      \(\overline{\{f(P)\}}\) is an irreducible component of \(X\);
\item every irreducible component \(Y\) of \(X\) through \(x\)
      is \(\overline{\{f(P)\}}\) for some minimal prime \(P\) of \(R\).
\end{enumerate}
This is the full original \(\mathrm{ChartOfFlat}\).  Neither
\(C\), \(R\), nor \(X\) is assumed Noetherian in this claim.

\begin{proof}
We recall explicitly the going-down consequence of flatness.
If \(P\in\operatorname{Spec}R\), with contraction \(p\),
and \(q\subset p\), there is \(P'\subset P\) contracting to
\(q\).  Indeed, \(C_p\to R_P\) is flat and local, hence
faithfully flat.  To see faithfulness, for any proper ideal
\(J\subset C_p\) its extension is contained in the maximal
ideal of \(R_P\).  A nonzero module contains a nonzero cyclic
submodule; flatness preserves that injection, and the cyclic
submodule remains nonzero after tensoring.  The fiber over
\(qC_p\) is consequently nonzero.  A prime of that fiber
provides the required \(P'\).

We also use the following elementary fact about the open chart.
If an irreducible set \(Y\subset X\) contains \(j(p)\), then
there is \(q\subset p\) in \(\operatorname{Spec}C\) such that
\(Y\subset\overline{\{j(q)\}}\).
The nonempty intersection \(Y\cap j(\operatorname{Spec}C)\)
is dense in \(Y\).  Its inverse image is irreducible; its
closure in \(\operatorname{Spec}C\) is \(V(q)\) for a prime
\(q\).  It contains \(p\), so \(q\subset p\).
Continuity of \(j\) and the stated density give the assertion.

Now let \(P\) be a minimal prime of \(R\), and suppose an
irreducible set \(Y\) contains \(\overline{\{f(P)\}}\).
Apply the preceding fact to obtain \(q\subset P\cap C\)
with \(Y\subset\overline{\{j(q)\}}\).  Going down gives
\(P'\subset P\) contracting to \(q\).  Minimality forces
\(P'=P\), so \(Y\subset\overline{\{f(P)\}}\).
The latter closed irreducible set is therefore maximal among
irreducible sets, proving the first property.

Conversely let \(Y\) be an irreducible component through \(x\).
The open-chart fact gives a prime
\(q\subset\mathfrak m_R\cap C\) with
\(Y\subset\overline{\{j(q)\}}\).  Going down from
\(\mathfrak m_R\) gives a prime \(P'\) over \(q\).
Choose a minimal prime \(P\subset P'\).
Specialization then gives
\(Y\subset\overline{\{f(P)\}}\).
Both maximality of \(Y\) and irreducibility of the latter
closure force equality.  Continuity and the closed-point
identity follow directly from the definition of \(f\).
\end{proof}
% END ACP16

% BEGIN ACP17
\paragraph{ACP17. The full transfer from a local chart to global components.}
\label{claim:ACP17}
Let \(T\subset X\) be open and \(G\subset X\) closed.
Let \(R\) be Noetherian and local and let
\(f:\operatorname{Spec}R\to X\) be any chart with all four
properties in ACP16, centered at \(x\).  Give \(R\) local
propagation data with type-I set \(f^{-1}(T)\), good set
\(f^{-1}(G)\), and threshold \(D\), and suppose the conclusion
of ACP14 holds for those data.  Then all three assertions hold:
\begin{enumerate}
\item If \(\operatorname{qdim}(P)\ge D\) for every minimal
      prime \(P\) of \(R\), goodness propagates between any
      two relevant components of \(X\) through \(x\).
\item If every minimal prime of \(R\) is type I, there is
      a relevant component of \(X\) through \(x\).
\item If there is a good type-I prime \(I\) of \(R\) with
      \(\operatorname{qdim}(I)\ge D\), there is a relevant
      component of \(X\) wholly contained in \(G\).
\end{enumerate}
Here relevant means meeting \(T\).  This is the full original
\(\mathrm{GlobOfCharts}\), for an arbitrary supplied chart;
it is not restricted to the flat charts just constructed.

\begin{proof}
An open set meeting \(\overline{\{y\}}\) contains \(y\).
Also, if \(y\in G\), closedness gives
\(\overline{\{y\}}\subset G\).
These two elementary observations account for both type and
goodness when passing between generic points and components.

For assertion 1, choose minimal primes \(P,P'\) with
\(Y=\overline{\{f(P)\}}\) and
\(Y'=\overline{\{f(P')\}}\), using the chart property for
components through \(x\).  Relevance and openness of \(T\)
make both primes type I.  If \(Y\subset G\), then \(P\)
is good, and its assumed dimension bound allows the local
propagation conclusion to be applied with seed \(P\).
It makes \(P'\) good.  Closedness of \(G\) then gives
\(Y'\subset G\).

For assertion 2, take an irreducible component through \(x\)
and use the chart property to write it as
\(\overline{\{f(P)\}}\) for a minimal \(P\).
By hypothesis \(f(P)\in T\), so this component is relevant.

For assertion 3, choose a minimal prime \(P\subset I\).
The inverse image of \(T\) is open, so it is stable under
generization; hence \(P\) is type I.  The local conclusion,
applied to the given seed \(I\), makes \(P\) good.
The chart's minimal-prime property makes
\(\overline{\{f(P)\}}\) an irreducible component.
It meets \(T\) at \(f(P)\) and is contained in \(G\),
proving the final assertion.
\end{proof}
% END ACP17

% BEGIN ACP18
\paragraph{ACP18. Combining the two rank-three loci gives one finite ideal.}
\label{claim:ACP18}
Let \(\mathcal B\) be any commutative ring with ideals
\(\mathfrak m,\lambda,M,I\).  If
\[
 \mathfrak m\subset\sqrt{M+\lambda},\qquad
 \mathfrak m\subset\sqrt{I+\lambda},
\]
then
\[
                  \mathfrak m\subset\sqrt{(M\cap I)+\lambda}.
                                                        \tag{ACP.26}
\]
In an original rank-three space let
\(\mathfrak e\subset\mathcal A\) be generated by the entries
whose weights are odd.  A type-I prime is off triangular and
does not contain \(\mathfrak e\); an exceptional-type prime
is off triangular and contains \(\mathfrak e\).
If \(M\) controls all type-I primes, \(I\) controls all
exceptional-type primes, and both satisfy \(\mathrm{FinLam}\),
then the single ideal \(M\cap I\) satisfies the full
\(\mathrm{OffFin}\) assertion (ACP.6).

\begin{proof}
Let \(P\) be any prime containing \((M\cap I)+\lambda\).
Since \(MI\subset M\cap I\), primality implies
\(M\subset P\) or \(I\subset P\).  Together with
\(\lambda\subset P\), the corresponding hypothesis forces
\(\mathfrak m\subset P\).  Intersect over all such primes
to obtain (ACP.26).

Every off-triangular prime of \(\mathcal A\) lies in exactly
one of the two indicated cases, according as it contains
\(\mathfrak e\).  In the first case its contraction contains
\(M\), in the second it contains \(I\); either contraction
therefore contains \(M\cap I\).  Equation (ACP.26) supplies
the remaining \(\mathrm{FinLam}\) condition on this same
ideal.  This is the actual final ideal operation in the
rank-three propagation theorem.  It does not by itself
construct the exceptional-type ideal \(I\).
\end{proof}
% END ACP18

\subsection{The exact arithmetic interface and the remaining geometric suppliers}

% BEGIN ACP19
\paragraph{ACP19. Forgetting the extra Alt properties preserves the common marking.}
\label{claim:ACP19}
At the original quantifiers,
\[
 \mathrm{GdEndAlt}\Longrightarrow\mathrm{GdEnd},
 \qquad \mathrm{HostEndAlt}\Longrightarrow\mathrm{HostEnd}.
                                                        \tag{ACP.27}
\]
These implications retain the identical \(g_1\), Marks and
FinEnd.  They do not supply either Alt endpoint.

\begin{proof}
The original rank-three automorphy predicate is
\(\mathrm{AutCsd}_0(\iota,r)\), the existence of the specified
regular algebraic conjugate-self-dual automorphic representation
attached to \(r\).  Its Alt version is literally
\[
 \mathrm{AutCsd}_0(\iota,r)\ \land
             \bigl(\mathrm{Continuous}(r)\Rightarrow
                                      \mathrm{WitA}(\iota,r)\bigr).
                                                        \tag{ACP.28}
\]
The premises of the original and Alt rank-three end conditions
are identical.  Projecting the first conjunct at each allowed
ordinary point gives the original rank-three end condition.
The Alt rank-two end condition likewise has the identical
premises and concludes the original attached-newform assertion
together with \(\mathrm{WitAQ}(\iota,W)\).  Project its first
conjunct at each allowed point and corner representation.

Now fix any initial datum under the original GdEnd quantifiers.
Use its Alt witness \(g_1\), retain Marks and FinEnd, and use
the two projections just described.  This proves GdEnd without
choosing another marker.  The Host statement is the same
argument with its rank-three conclusion only.  Nothing in
this projection constructs the extra \(\mathrm{WitA}\) or
\(\mathrm{WitAQ}\), and neither converse is asserted.
The actual terminal route continues to retain those properties.
\end{proof}
% END ACP19

\paragraph{Where the proved propagation principles are used.}
For the arithmetic spaces, the global topological space is
\(X=\operatorname{Spec}\mathcal A\), the good locus is
\(V(M\mathcal A)\), and the core is
\[
 \{\mathfrak p:\mathfrak m_{\mathcal B}
        \subset\pi(\mathfrak p),\quad
                      \mathfrak t_+\not\subset\mathfrak p\}.
                                                        \tag{ACP.29}
\]
The chosen points are its maximal ideals.  In rank two the
exceptional locus is empty and the local threshold is
\(2d+1\); in rank three the exceptional locus is
\(V(\mathfrak e)\) and the threshold is \(3d+1\).
The original \(P1'\) and \(P23'\) conditions supply the
two local propagation rules in ACP14.  ACP13 supplies the
prime-avoidance mechanism once the required bounds on the
forbidden loci have been proved.  ACP16--17 transfer that
local result to the component-incidence property used by
ACP15.  ACP18 is the final combination of the type-I and
exceptional-type ideals.  ACP11 then supplies the exact
finite-point conclusion used by AFE20.

The still needed geometric inputs are specific.  They are:
the representable local charts and their flatness; the sharp
component-chain and dimension estimates; the bounds for the
triangular and coincident-character loci; connectedness and
the indicated closed-point density of the core; the passage
from the seed prime to a sufficiently large local prime; and
the actual rank-two-to-rank-three corner construction producing
the exceptional-type ideal.  They are the inputs proved in
the source files \texttt{LocFlat}, \texttt{ChainedBot},
\texttt{ChainedOfComap},
\texttt{KrullDimLtOfNotLeTriangular},
\texttt{PreconnectedRankTwo}, \texttt{Seedpt}, and
\texttt{QuotientHomNonvanishingBound}, respectively.
This chapter proves neither the full rank-two nor the full
rank-three propagation theorem without those constructions.

The actual later call keeps a single extended frame for both
ranks, while allowing different spaces and modularity ideals.
The rank-two propagation result feeds the rank-three corner
and propagation arguments.  Their \(\mathrm{OffFin}\) result
then enters ACP11.  The automorphic seeds, the modularity
realizations and arithmetic-point classicality remain separate
suppliers.  The completed Dat and Host frame entries, including
Host \(\ell=3\), are reused; no new restriction on their
marked element, specified place, or order of choices has
been introduced.

